% ============================================================================
% Related Work (rewritten): survey the landscape as prose first; position OUR
% contribution only in the final paragraph. Drop into Overleaf and either
%   % ============================================================================
% Related Work (rewritten): survey the landscape as prose first; position OUR
% contribution only in the final paragraph. Drop into Overleaf and either
%   % ============================================================================
% Related Work (rewritten): survey the landscape as prose first; position OUR
% contribution only in the final paragraph. Drop into Overleaf and either
%   % ============================================================================
% Related Work (rewritten): survey the landscape as prose first; position OUR
% contribution only in the final paragraph. Drop into Overleaf and either
%   \input{relatedwork}
% in place of the old "\subsection{Related work} ... \end{table}" block, or
% paste the body below over it. Uses only refs already in refs.bib.
% ============================================================================
\subsection{Related work}
\label{sec:related}
Automated game testing has traditionally leaned on hand-written oracles: scripted assertions and record-and-replay
regression suites that re-run fixed trajectories, invariant and property checks, and fuzzing that perturbs inputs to
force a crash, all backed by human QA playthroughs. Such methods catch what their authors anticipated, but none
\emph{explores} the state space on its own, so the rare states that trigger a bug, or a level's own win condition, are
reached only when a person or a script already knew to go there. The open problem is therefore an exploration policy that
drives a game into those rare states within the small interaction budget a testing pass can afford, and recent work
attacks it by driving the game with an agent, including per-step large-language-model agents that act as
testers~\citep{titan}.

The dominant learned approach trains a reinforcement-learning explorer whose reward is augmented with an
intrinsic-motivation bonus for novelty: count-based exploration~\citep{tang2017}, random network
distillation~\citep{rnd}, and curiosity from prediction error~\citep{icm}. These are effective but sample-hungry,
typically needing $10^6$--$10^8$ environment steps before the bonus becomes informative, far more interaction than a
testing pass allocates. A complementary line stores and revisits promising states instead of learning a policy:
Go-Explore~\citep{goexplore} keeps an archive of returnable states and is the strongest classical exploration family,
though its results hinge on a hand-chosen cell abstraction and are reported per game. Simpler still, action-repeat
(sticky actions), introduced as a stochasticity control for the Arcade Learning Environment~\citep{machado2018}, is not
an exploration method at all, yet is a surprisingly strong baseline on physics-like dynamics.

A separate thread has large language models \emph{write code} for control rather than act step by step.
FunSearch~\citep{funsearch} evolves programs against an executable score; Eureka~\citep{eureka} writes reward functions
that are then optimized by standard RL; and PromptPO~\citep{promptpo} asks when an LLM can itself serve as the policy
optimizer for sequential RL tasks. Together these show that an LLM-authored artifact can rival hand-designed or trained
controllers, but they target return maximization or reward design, and the resulting policy is still trained or the LLM
is still in the loop at inference time.

Our setting sits at a point none of these occupy. We ask for an exploration procedure that is evaluated at a fixed
interaction budget, aimed at reaching sparse target and defect states rather than maximizing return, written once as a
program and run \emph{training-free}, and execution-verified, with every counted hit replayed from a clean reset.
Table~\ref{tab:caps} places the families above on exactly these axes; the region that is budget-matched, training-free,
target-reaching, and replay-verified at once is the one our written tester occupies.

\begin{table}[t]\centering\small
\caption{Where each family sits. Our tester is the one that is budget-matched, training-free, and target-reaching at
once; universality across all tasks is \emph{not} claimed (Section~\ref{sec:result}).}
\label{tab:caps}
\setlength{\tabcolsep}{4pt}
\resizebox{\linewidth}{!}{%
\begin{tabular}{@{}lcccc@{}}\toprule
 & no per-env & fixed interaction & sparse-reward & replay-\\
 & training & budget & target-reaching & verified \\ \midrule
PPO + intrinsic~\citep{tang2017,rnd,icm} & \ding{55} & \ding{55} & partial & \ding{55} \\
Go-Explore~\citep{goexplore} & \checkmark & partial & \checkmark & \ding{55} \\
Action-repeat~\citep{machado2018} & \checkmark & \checkmark & \ding{55} & \ding{55} \\
Eureka / PromptPO~\citep{eureka,promptpo} & \ding{55}/\checkmark & \ding{55} & \ding{55} & \ding{55} \\
\textbf{Written program (ours)} & \checkmark & \checkmark & \checkmark & \checkmark \\ \bottomrule
\end{tabular}}
\end{table}

% in place of the old "\subsection{Related work} ... \end{table}" block, or
% paste the body below over it. Uses only refs already in refs.bib.
% ============================================================================
\subsection{Related work}
\label{sec:related}
Automated game testing has traditionally leaned on hand-written oracles: scripted assertions and record-and-replay
regression suites that re-run fixed trajectories, invariant and property checks, and fuzzing that perturbs inputs to
force a crash, all backed by human QA playthroughs. Such methods catch what their authors anticipated, but none
\emph{explores} the state space on its own, so the rare states that trigger a bug, or a level's own win condition, are
reached only when a person or a script already knew to go there. The open problem is therefore an exploration policy that
drives a game into those rare states within the small interaction budget a testing pass can afford, and recent work
attacks it by driving the game with an agent, including per-step large-language-model agents that act as
testers~\citep{titan}.

The dominant learned approach trains a reinforcement-learning explorer whose reward is augmented with an
intrinsic-motivation bonus for novelty: count-based exploration~\citep{tang2017}, random network
distillation~\citep{rnd}, and curiosity from prediction error~\citep{icm}. These are effective but sample-hungry,
typically needing $10^6$--$10^8$ environment steps before the bonus becomes informative, far more interaction than a
testing pass allocates. A complementary line stores and revisits promising states instead of learning a policy:
Go-Explore~\citep{goexplore} keeps an archive of returnable states and is the strongest classical exploration family,
though its results hinge on a hand-chosen cell abstraction and are reported per game. Simpler still, action-repeat
(sticky actions), introduced as a stochasticity control for the Arcade Learning Environment~\citep{machado2018}, is not
an exploration method at all, yet is a surprisingly strong baseline on physics-like dynamics.

A separate thread has large language models \emph{write code} for control rather than act step by step.
FunSearch~\citep{funsearch} evolves programs against an executable score; Eureka~\citep{eureka} writes reward functions
that are then optimized by standard RL; and PromptPO~\citep{promptpo} asks when an LLM can itself serve as the policy
optimizer for sequential RL tasks. Together these show that an LLM-authored artifact can rival hand-designed or trained
controllers, but they target return maximization or reward design, and the resulting policy is still trained or the LLM
is still in the loop at inference time.

Our setting sits at a point none of these occupy. We ask for an exploration procedure that is evaluated at a fixed
interaction budget, aimed at reaching sparse target and defect states rather than maximizing return, written once as a
program and run \emph{training-free}, and execution-verified, with every counted hit replayed from a clean reset.
Table~\ref{tab:caps} places the families above on exactly these axes; the region that is budget-matched, training-free,
target-reaching, and replay-verified at once is the one our written tester occupies.

\begin{table}[t]\centering\small
\caption{Where each family sits. Our tester is the one that is budget-matched, training-free, and target-reaching at
once; universality across all tasks is \emph{not} claimed (Section~\ref{sec:result}).}
\label{tab:caps}
\setlength{\tabcolsep}{4pt}
\resizebox{\linewidth}{!}{%
\begin{tabular}{@{}lcccc@{}}\toprule
 & no per-env & fixed interaction & sparse-reward & replay-\\
 & training & budget & target-reaching & verified \\ \midrule
PPO + intrinsic~\citep{tang2017,rnd,icm} & \ding{55} & \ding{55} & partial & \ding{55} \\
Go-Explore~\citep{goexplore} & \checkmark & partial & \checkmark & \ding{55} \\
Action-repeat~\citep{machado2018} & \checkmark & \checkmark & \ding{55} & \ding{55} \\
Eureka / PromptPO~\citep{eureka,promptpo} & \ding{55}/\checkmark & \ding{55} & \ding{55} & \ding{55} \\
\textbf{Written program (ours)} & \checkmark & \checkmark & \checkmark & \checkmark \\ \bottomrule
\end{tabular}}
\end{table}

% in place of the old "\subsection{Related work} ... \end{table}" block, or
% paste the body below over it. Uses only refs already in refs.bib.
% ============================================================================
\subsection{Related work}
\label{sec:related}
Automated game testing has traditionally leaned on hand-written oracles: scripted assertions and record-and-replay
regression suites that re-run fixed trajectories, invariant and property checks, and fuzzing that perturbs inputs to
force a crash, all backed by human QA playthroughs. Such methods catch what their authors anticipated, but none
\emph{explores} the state space on its own, so the rare states that trigger a bug, or a level's own win condition, are
reached only when a person or a script already knew to go there. The open problem is therefore an exploration policy that
drives a game into those rare states within the small interaction budget a testing pass can afford, and recent work
attacks it by driving the game with an agent, including per-step large-language-model agents that act as
testers~\citep{titan}.

The dominant learned approach trains a reinforcement-learning explorer whose reward is augmented with an
intrinsic-motivation bonus for novelty: count-based exploration~\citep{tang2017}, random network
distillation~\citep{rnd}, and curiosity from prediction error~\citep{icm}. These are effective but sample-hungry,
typically needing $10^6$--$10^8$ environment steps before the bonus becomes informative, far more interaction than a
testing pass allocates. A complementary line stores and revisits promising states instead of learning a policy:
Go-Explore~\citep{goexplore} keeps an archive of returnable states and is the strongest classical exploration family,
though its results hinge on a hand-chosen cell abstraction and are reported per game. Simpler still, action-repeat
(sticky actions), introduced as a stochasticity control for the Arcade Learning Environment~\citep{machado2018}, is not
an exploration method at all, yet is a surprisingly strong baseline on physics-like dynamics.

A separate thread has large language models \emph{write code} for control rather than act step by step.
FunSearch~\citep{funsearch} evolves programs against an executable score; Eureka~\citep{eureka} writes reward functions
that are then optimized by standard RL; and PromptPO~\citep{promptpo} asks when an LLM can itself serve as the policy
optimizer for sequential RL tasks. Together these show that an LLM-authored artifact can rival hand-designed or trained
controllers, but they target return maximization or reward design, and the resulting policy is still trained or the LLM
is still in the loop at inference time.

Our setting sits at a point none of these occupy. We ask for an exploration procedure that is evaluated at a fixed
interaction budget, aimed at reaching sparse target and defect states rather than maximizing return, written once as a
program and run \emph{training-free}, and execution-verified, with every counted hit replayed from a clean reset.
Table~\ref{tab:caps} places the families above on exactly these axes; the region that is budget-matched, training-free,
target-reaching, and replay-verified at once is the one our written tester occupies.

\begin{table}[t]\centering\small
\caption{Where each family sits. Our tester is the one that is budget-matched, training-free, and target-reaching at
once; universality across all tasks is \emph{not} claimed (Section~\ref{sec:result}).}
\label{tab:caps}
\setlength{\tabcolsep}{4pt}
\resizebox{\linewidth}{!}{%
\begin{tabular}{@{}lcccc@{}}\toprule
 & no per-env & fixed interaction & sparse-reward & replay-\\
 & training & budget & target-reaching & verified \\ \midrule
PPO + intrinsic~\citep{tang2017,rnd,icm} & \ding{55} & \ding{55} & partial & \ding{55} \\
Go-Explore~\citep{goexplore} & \checkmark & partial & \checkmark & \ding{55} \\
Action-repeat~\citep{machado2018} & \checkmark & \checkmark & \ding{55} & \ding{55} \\
Eureka / PromptPO~\citep{eureka,promptpo} & \ding{55}/\checkmark & \ding{55} & \ding{55} & \ding{55} \\
\textbf{Written program (ours)} & \checkmark & \checkmark & \checkmark & \checkmark \\ \bottomrule
\end{tabular}}
\end{table}

% in place of the old "\subsection{Related work} ... \end{table}" block, or
% paste the body below over it. Uses only refs already in refs.bib.
% ============================================================================
\subsection{Related work}
\label{sec:related}
Automated game testing has traditionally leaned on hand-written oracles: scripted assertions and record-and-replay
regression suites that re-run fixed trajectories, invariant and property checks, and fuzzing that perturbs inputs to
force a crash, all backed by human QA playthroughs. Such methods catch what their authors anticipated, but none
\emph{explores} the state space on its own, so the rare states that trigger a bug, or a level's own win condition, are
reached only when a person or a script already knew to go there. The open problem is therefore an exploration policy that
drives a game into those rare states within the small interaction budget a testing pass can afford, and recent work
attacks it by driving the game with an agent, including per-step large-language-model agents that act as
testers~\citep{titan}.

The dominant learned approach trains a reinforcement-learning explorer whose reward is augmented with an
intrinsic-motivation bonus for novelty: count-based exploration~\citep{tang2017}, random network
distillation~\citep{rnd}, and curiosity from prediction error~\citep{icm}. These are effective but sample-hungry,
typically needing $10^6$--$10^8$ environment steps before the bonus becomes informative, far more interaction than a
testing pass allocates. A complementary line stores and revisits promising states instead of learning a policy:
Go-Explore~\citep{goexplore} keeps an archive of returnable states and is the strongest classical exploration family,
though its results hinge on a hand-chosen cell abstraction and are reported per game. Simpler still, action-repeat
(sticky actions), introduced as a stochasticity control for the Arcade Learning Environment~\citep{machado2018}, is not
an exploration method at all, yet is a surprisingly strong baseline on physics-like dynamics.

A separate thread has large language models \emph{write code} for control rather than act step by step.
FunSearch~\citep{funsearch} evolves programs against an executable score; Eureka~\citep{eureka} writes reward functions
that are then optimized by standard RL; and PromptPO~\citep{promptpo} asks when an LLM can itself serve as the policy
optimizer for sequential RL tasks. Together these show that an LLM-authored artifact can rival hand-designed or trained
controllers, but they target return maximization or reward design, and the resulting policy is still trained or the LLM
is still in the loop at inference time.

Our setting sits at a point none of these occupy. We ask for an exploration procedure that is evaluated at a fixed
interaction budget, aimed at reaching sparse target and defect states rather than maximizing return, written once as a
program and run \emph{training-free}, and execution-verified, with every counted hit replayed from a clean reset.
Table~\ref{tab:caps} places the families above on exactly these axes; the region that is budget-matched, training-free,
target-reaching, and replay-verified at once is the one our written tester occupies.

\begin{table}[t]\centering\small
\caption{Where each family sits. Our tester is the one that is budget-matched, training-free, and target-reaching at
once; universality across all tasks is \emph{not} claimed (Section~\ref{sec:result}).}
\label{tab:caps}
\setlength{\tabcolsep}{4pt}
\resizebox{\linewidth}{!}{%
\begin{tabular}{@{}lcccc@{}}\toprule
 & no per-env & fixed interaction & sparse-reward & replay-\\
 & training & budget & target-reaching & verified \\ \midrule
PPO + intrinsic~\citep{tang2017,rnd,icm} & \ding{55} & \ding{55} & partial & \ding{55} \\
Go-Explore~\citep{goexplore} & \checkmark & partial & \checkmark & \ding{55} \\
Action-repeat~\citep{machado2018} & \checkmark & \checkmark & \ding{55} & \ding{55} \\
Eureka / PromptPO~\citep{eureka,promptpo} & \ding{55}/\checkmark & \ding{55} & \ding{55} & \ding{55} \\
\textbf{Written program (ours)} & \checkmark & \checkmark & \checkmark & \checkmark \\ \bottomrule
\end{tabular}}
\end{table}
